
\documentclass{article}
\usepackage{colortbl}
\usepackage{makecell}
\usepackage{multirow}
\usepackage{supertabular}

\begin{document}

\newcounter{utterance}

\centering \large Interaction Transcript for game `cladder', experiment `full\_v1.5\_default', episode 5136 with Qwen3.5{-}9B.
\vspace{24pt}

{ \footnotesize  \setcounter{utterance}{1}
\setlength{\tabcolsep}{0pt}
\begin{supertabular}{c@{$\;$}|p{.15\linewidth}@{}p{.15\linewidth}p{.15\linewidth}p{.15\linewidth}p{.15\linewidth}p{.15\linewidth}}
    \# & \multicolumn{2}{c}{Player} && \multicolumn{2}{c}{Game Master} \\
    \hline

    \theutterance \stepcounter{utterance}  
    & & & \multicolumn{4}{p{0.6\linewidth}}{
        \cellcolor[rgb]{0.9,0.9,0.9}{
            \makecell[{{p{\linewidth}}}]{
                \texttt{\tiny{[P1$\langle$GM]}}
                \texttt{You are an expert in causal inference. The following question is not a typical commonsense query, but rather a meticulously designed question created by a professor specializing in causal inference, intended to assess the students' mastery of the course content.} \\
\texttt{Imagine a self{-}contained, hypothetical world with only the following conditions, and without any unmentioned factors or causal relationships: Unobserved confounders has a direct effect on education level and foot size. Proximity to a college has a direct effect on education level. Education level has a direct effect on foot size. Unobserved confounders is unobserved. The overall probability of college degree or higher is 47\%. The probability of high school degree or lower and large feet is 32\%. The probability of college degree or higher and large feet is 10\%. Is the chance of large feet smaller when observing college degree or higher?} \\
\texttt{Start your answer with 'Yes' or 'No', followed by additional reasoning or evidence to support your explanation.} \\
            }
        }
    }
     \\ \\

    \theutterance \stepcounter{utterance}  
    & \multicolumn{4}{p{0.6\linewidth}}{
        \cellcolor[rgb]{0.561,0.737,0.561}{
            \makecell[{{p{\linewidth}}}]{
                \texttt{\tiny{[P1$\rangle$GM]}}
                \texttt{No} \\
\\ 
\texttt{**Reasoning:**} \\
\\ 
\texttt{To determine if the chance of large feet is smaller when observing a college degree or higher, we must compare the conditional probability \$P(\textbackslash{}text\{Large Feet\} \textbackslash{}mid \textbackslash{}text\{College\})\$ with the marginal probability \$P(\textbackslash{}text\{Large Feet\})\$.} \\
\\ 
\texttt{1.  **Calculate the Marginal Probability of Large Feet (\$P(\textbackslash{}text\{Large Feet\})\$):**} \\
\texttt{    The problem provides the probability of the union of two mutually exclusive events regarding foot size (assuming "large feet" and "not large feet" partition the space, or at least that the provided probabilities cover the relevant cases):} \\
\texttt{    *   \$P(\textbackslash{}text\{High School or Lower\} \textbackslash{}cap \textbackslash{}text\{Large Feet\}) = 0.32\$} \\
\texttt{    *   \$P(\textbackslash{}text\{College or Higher\} \textbackslash{}cap \textbackslash{}text\{Large Feet\}) = 0.10\$} \\
\\ 
\texttt{    Assuming these are the only two categories for education level mentioned in the context of the foot size data (or that "High School or Lower" and "College or Higher" are the exhaustive groups provided for this calculation), the total probability of having large feet is the sum of these disjoint probabilities:} \\
\texttt{    \$\$P(\textbackslash{}text\{Large Feet\}) = 0.32 + 0.10 = 0.42\$\$} \\
\\ 
\texttt{2.  **Calculate the Conditional Probability of Large Feet given College (\$P(\textbackslash{}text\{Large Feet\} \textbackslash{}mid} \\
            }
        }
    }
    & & \\ \\

    \theutterance \stepcounter{utterance}  
    & & & \multicolumn{2}{p{0.3\linewidth}}{
        \cellcolor[rgb]{0.9,0.9,0.9}{
            \makecell[{{p{\linewidth}}}]{
                \texttt{\tiny{[GM$|$GM]}}
                \texttt{No} \\
            }
        }
    }
    & & \\ \\

    \theutterance \stepcounter{utterance}  
    & & & \multicolumn{2}{p{0.3\linewidth}}{
        \cellcolor[rgb]{0.9,0.9,0.9}{
            \makecell[{{p{\linewidth}}}]{
                \texttt{\tiny{[GM$|$GM]}}
                \texttt{yes} \\
            }
        }
    }
    & & \\ \\

    \theutterance \stepcounter{utterance}  
    & & & \multicolumn{2}{p{0.3\linewidth}}{
        \cellcolor[rgb]{0.9,0.9,0.9}{
            \makecell[{{p{\linewidth}}}]{
                \texttt{\tiny{[GM$|$GM]}}
                \texttt{game\_result = LOSE} \\
            }
        }
    }
    & & \\ \\

\end{supertabular}
}

\end{document}
